\chapter{Introdução}

\section{Contextualização}

A organização de festivais envolve a coordenação de um elevado volume de
informação: datas e locais, palcos, artistas e respetiva programação,
inscrições, pagamentos e o controlo de acesso ao recinto. Quando esta
informação é gerida de forma dispersa (folhas de cálculo, mensagens, listas em
papel), surgem erros frequentes, como inscrições duplicadas, bilhetes
reutilizados ou programações com conflitos de horário.

O presente trabalho propõe um \textbf{Sistema de Gestão de Festivais}: uma
aplicação \emph{web} que centraliza esta informação numa base de dados
relacional e disponibiliza, a cada tipo de utilizador, apenas as
funcionalidades que lhe dizem respeito. O público consulta festivais e
programação; os participantes inscrevem-se, efetuam um pagamento simulado e
recebem um bilhete com código QR; os organizadores gerem os seus festivais; o
\emph{Staff} valida as entradas no recinto; e o administrador supervisiona a
plataforma.

\section{Objetivos}

Os objetivos do trabalho são:
\begin{itemize}
  \item modelar o sistema de informação, identificando atores, requisitos,
        casos de uso e respetivos cenários;
  \item conceber e implementar uma base de dados relacional que garanta a
        integridade dos dados (inscrições únicas, bilhetes de utilização
        única, programação sem conflitos);
  \item desenhar a interface gráfica (\emph{wireframes}, formatos de saída e
        \emph{wireflow});
  \item implementar uma aplicação \emph{web} responsiva e multilingue
        (Português e Inglês), com autenticação local e controlo de acesso por
        perfil;
  \item gerir o projeto de forma colaborativa, com controlo de versões e um
        quadro de tarefas.
\end{itemize}

\section{Âmbito}

Ficam dentro do âmbito: a gestão de festivais, palcos, artistas e programação;
o registo e autenticação de utilizadores; as inscrições com pagamento
\textbf{simulado}; a emissão de bilhetes com código QR; a validação de entradas;
e a administração de contas e perfis.

Ficam fora do âmbito: pagamentos reais e reembolsos, venda de diferentes tipos
de bilhete (por exemplo, diário ou VIP) e integração com sistemas externos de
bilhética.

\section{Metodologia e Gestão do Projeto}

\pendente{Descrever a abordagem iterativa adotada (Scrum/Kanban), o
\emph{Product Backlog} com as histórias de utilizador e critérios de aceitação,
a distribuição de tarefas pelos elementos do grupo e as ferramentas usadas
(GitHub, Jira ou quadro Kanban). Incluir uma captura do quadro.}
